\documentclass[compacto,10pt]{pubciencias-notas}

% -- Paquetes adicionales
\usepackage{comandosEPN}

% -- Datos del libro
\autor{Departamento de Formación Básica}
\materia{Álgebra Lineal}
\nota{Semana 14 Clase 01}
\fecha{Semestre 2024-B}

\logouno[1.1cm]{Logo_EPN.eps}
\logodos[1.7cm]{Logo_DFB.eps}

\definecolor{colordef}{cmyk}{0.81,0.62,0.00,0.22}

\begin{document}

\encabezado
\vspace*{-8mm}

\setcounter{section}{6}
\section{Valores y vectores propios}


\setcounter{subsection}{1}
\subsection{Polinomio caracter\'istico}


\begin{defi}
Sea $A\in \Mat[\K]{n}{n}$. Al polinomio
\[
    p(\lambda) = \det(A - \lambda I_{n})
\]
se lo denomina polinomio característico de $A$, donde $I_{n}$ es la matriz identidad.
\end{defi}

\begin{ejem}\hspace{0pt} Dada la matriz $A = \begin{pmatrix} 1&2\\4&3 \end{pmatrix}$, su polinomio característico es:
\begin{align*}
    p(\lambda)  &=  \det(A - \lambda I_{2})\\
                &=  \det\l(\begin{pmatrix}1&2\\4&3\end{pmatrix}-\lambda\begin{pmatrix}1&0\\0&1\end{pmatrix}\r)\\
                &=  \det\begin{pmatrix}1-\lambda &2 \\ 4 & 3-\lambda \end{pmatrix}\\
                &= \lambda^2 - 4\lambda - 5.\\
%                &= (\lambda + 1) (\lambda - 5).
\end{align*}
Por lo tanto, $p(\lambda)= \lambda^2 - 4\lambda - 5$ es el polinomio característico de $A$.
\end{ejem}


\begin{prop}
Sean $A\in \Mat[\K]{n}{n}$ y $\lambda \in \K$. Se tiene que $\lambda$ es un valor propio de $A$ si y solo si
\[
    \det(A - \lambda I_{n}) = 0,
\]
donde $I_{n}$ es la matriz identidad.
\end{prop}

\begin{advertencia}
Notemos que todo valor propio de $A\in\Mat[\K]{n}{n}$, es raíz del polinomio característico de $A$.
\end{advertencia}

\begin{ejem}\hspace{0pt} Dada la matriz $A = \begin{pmatrix} 1&2\\4&3 \end{pmatrix}$, del ejemplo anterior su polinomio característico es:
\begin{align*}
    p(\lambda)  &=  \det(A - \lambda I_{2})\\
%                &=  \det\l(\begin{pmatrix}1&2\\4&3\end{pmatrix}-\lambda\begin{pmatrix}1&0\\0&1\end{pmatrix}\r)\\
%                &=  \det\begin{pmatrix}1-\lambda &2 \\ 4 & 3-\lambda \end{pmatrix}\\
                &= \lambda^2 - 4\lambda - 5\\
                &= (\lambda + 1) (\lambda - 5).
\end{align*}
Por lo tanto, $\lambda_1 = -1$ y $\lambda_2 = 5$ son los valores propios de $A$.
\end{ejem}

\begin{prop}
Los valores propios de una matriz triangular son los elementos de la diagonal principal de la matriz. 
\end{prop}
%
%\begin{teo}[Teorema de Cayley–Hamilton]
%Sea $A\in \Mat[\K]{n}{n}$ y $p(\lambda)$ su polinomio característico. Se tiene que
%\[
 %   p(A) = 0.
%\]
%\end{teo}

%\begin{advertencia}
%Es decir, toda matriz es un cero de su polinomio característico, esto es,
%si $A\in\Mat[\K]{n}{n}$ y si su polinomio característico es
%\[
%   p(\lambda) =  \lambda^{n} + c_{n-1}\lambda^{n-1} + \cdots + c_{1}\lambda + c_{0}
%\]
%entonces,
%\[
%   p(A) =  A^{n} + c_{n-1}A^{n-1} + \cdots + c_{1}A + c_{0}I_{n} = 0.
%\]
%\end{advertencia}

%\begin{ejem}\hspace{0pt} 
%Para la matriz $A = \begin{pmatrix} 1&2\\4&3 \end{pmatrix}$, sabemos que su polinomio característico es:
%\[
%    p(\lambda)  = \lambda^2 - 4\lambda - 5.
%\]
%Luego, se tiene que:
%\[
%    p(A)  = A^2 - 4A - 5I_{2} = \begin{pmatrix}9&8\\16&17\end{pmatrix} - 4\begin{pmatrix}1&2\\4&3\end{pmatrix} - 5\begin{pmatrix}1&0\\0&1\end{pmatrix} = \begin{pmatrix}0&0\\0&0\end{pmatrix}.
%\]
%\end{ejem}

\begin{advertencia}
\textbf{Algoritmo para encontrar los valores y vectores propios de una matriz.}

Sea $A\in\Mat[\K]{n}{n}$.
\begin{description}
    \item[Paso 1.] Hallar el polinomio característico $p(\lambda)$ de $A$.
    \item[Paso 2.] Hallar las raíces de $p(\lambda)$.
    \item[Paso 3.] Repetir los pasos $a)$ y $b)$ para cada valor propio $\lambda$ de $A$:
        \begin{enumerate}
            \item[$a)$] Construir $M = A - \lambda I_{n}$ restando $\lambda$ a los elementos diagonales de $A$.
            \item[$b)$] Encontrar una base para el espacio solución de $Mv = 0$. Los vectores de la base son vectores propios linealmente independientes de $A$ asociados a $\lambda$.
        \end{enumerate}
\end{description}
\end{advertencia}




\begin{ejem}\hspace{0pt} Dada la matriz $A = \begin{pmatrix}
1&-1&4 \\ 3&2&-1 \\ 2&1&-1 \end{pmatrix}$. 

Paso 1. Su polinomio característico es:
\begin{align*}
    p(\lambda)  &=  \det(A - \lambda I_{3})\\
                &=  \det\l(\begin{pmatrix}
1&-1&4 \\ 3&2&-1 \\ 2&1&-1 \end{pmatrix}-\lambda\begin{pmatrix} 1&0&0 \\ 0&1&0 \\ 0&0&1 \end{pmatrix}\r)\\
                &=  \det\begin{pmatrix} 1-\lambda &-1&4 \\ 3&2-\lambda &-1 \\ 2&1&-1-\lambda  \end{pmatrix}\\
                &= -\lambda^3+2\lambda^2 + 5\lambda - 6.\\
%                &= (\lambda + 1) (\lambda - 5).
\end{align*}
Por lo tanto, $p(\lambda)=-\lambda^3+2\lambda^2 + 5\lambda - 6$.

Paso 2. Las raíces de $p(\lambda)$ son:


\begin{align*}
    p(\lambda)  &=  0\\
%                &=  \det\l(\begin{pmatrix}1&2\\4&3\end{pmatrix}-\lambda\begin{pmatrix}1&0\\0&1\end{pmatrix}\r)\\
%                &=  \det\begin{pmatrix}1-\lambda &2 \\ 4 & 3-\lambda \end{pmatrix}\\
              -\lambda^3+2\lambda^2 + 5\lambda - 6 &=0 \\
               -(\lambda-1)(\lambda+2)(\lambda -3) &=0 .
\end{align*}




por lo tanto las raíces de $p(\lambda)$ son $\lambda_{1}=1,\lambda_{2}=-2$ y $\lambda_{3}=3$.


Paso 3.

Para $\lambda_{1}=1$.

\[
A-\lambda_{1}I_{3}=\begin{pmatrix}
1&-1&4 \\ 3&2&-1 \\ 2&1&-1 \end{pmatrix} -\begin{pmatrix}
1&0&0 \\ 0&1&0 \\ 0&0&1 \end{pmatrix}=\begin{pmatrix}
0&-1&4 \\ 3&1&-1 \\ 2&1&-2 \end{pmatrix}.
\]

Encontremos la solución de $(A-\lambda_{1}I_{3})v=0$

\[
(A-\lambda_{1}I_{3})v= \begin{pmatrix}
0&-1&4 \\ 3&1&-1 \\ 2&1&-2 \end{pmatrix}\begin{pmatrix} x_{1}\\ x_{2}\\ x_{3}\end{pmatrix}=\begin{pmatrix} 0\\ 0\\ 0\end{pmatrix}
\]

por lo que la matriz ampliada del sistema es


\begin{align*}
\begin{pmatrix}0 & -1 & 4 & | & 0\\ 3 & 1 & -1 & | & 0\\ 2 & 1 & -2 & | & 0\end{pmatrix}
& \sim 
\begin{pmatrix}0 & -1 & 4 & | & 0\\ 3 & 0 & 3 & | & 0\\ 2 & 0 & 2 & | & 0\end{pmatrix}
&& \begin{array}{c}
    F_1+F_2\rightarrow F_2\\
    F_1+F_3\rightarrow F_3
\end{array}\\
& \sim 
\begin{pmatrix}0 & -1 & 4 & | & 0\\ 1 & 0 & 1 & | & 0\\ 2 & 0 & 2 & | & 0\end{pmatrix}
&& \frac{1}{3}F_2\rightarrow F_2\\
& \sim 
\begin{pmatrix}0 & -1 & 4 & | & 0\\ 1 & 0 & 1 & | & 0\\ 0 & 0 & 0 & | & 0\end{pmatrix}
&& -2F_2+F_3\rightarrow F_3
\end{align*}

de donde se tiene que $x_{1}=-x_{3}$; $x_{2}=4x_{3}$ y
\[
    E_1 = \left \langle (-1,4,1) \right \rangle.
\]

Verifique que $ E_{-2} =\left  \langle (1,-1,-1)  \right \rangle$ y $ E_3=\left \langle (1,2,1)  \right \rangle.$



% y $ E_3 \right \rangle
% =\left \langle  \right \rangle gen\l\{(1,2,1)\r\}.$

En conclusión $\lambda_{1}=1,\lambda_{2}=-2$ y $\lambda_{3}=3$ son los valores propios de $A$ y $(-1,4,1),(1,-1,-1),(1,2,1)$ son vectores propios de
$\lambda_{1},\lambda_{2}$ y $\lambda_{3}$, respectivamente. 

\end{ejem}






\begin{ejer}\hspace{0pt}
Dada la aplicación lineal:
\[
    \funcion{T}{\R^2}{\R^2}{(a,b)}{(6a-b,3a+2b)}
\]
Hallar los valores y vectores propios de $T$.
\end{ejer}

\begin{teo}
Cualquier polinomio de grado $n$ tiene exactamente $n$ raíces (contando con las veces que una determinada raíz se repite) en los complejos.
\end{teo}

\begin{defi}
Sea $\lambda$ un valor propio de la matriz $A$, el número de veces que un valor propio $\lambda$ se repite como raíz del polinomio característico $\rho(\lambda)$ se llama multiplicidad algebraica, notada $m_a(\lambda)$.
\end{defi}

\begin{defi}
Sea $\lambda$ un valor propio de la matriz $A$, entonces la multiplicidad geométrica de $\lambda$, notada $m_g(\lambda)$ está dada por
\[ m_g(\lambda) = \dim(E_\lambda)\]
\end{defi}

\begin{teo}
Sea $\lambda$ valor propio de la matriz $A$, entonces
\begin{center}
 $m_g(\lambda) \leq m_a(\lambda)$
\end{center}
\end{teo}

\begin{teo}
Contando multiplicidades, toda matriz $A \in \Mat[\K]{n}{n} $ tiene exactamente $n$ valores propios en los complejos.
\end{teo}

\begin{ejem}\hspace{0pt}\label{Cap7:ejem03} La matriz $A = \begin{pmatrix}1&0&0&1\\0&2&0&0\\0&0&2&0\\1&0&0&1\end{pmatrix}$ tiene como valores propios a $\lambda_{1} = 0$ y $\lambda_{2}=2$ con $m_a(\lambda_{1})=1$, $m_a(\lambda_{2})=3$, $m_g(\lambda_{1})=1$ y $m_g(\lambda_{2})=3$.
\end{ejem}


\end{document}